\documentclass[UTF8,fontset=fandol,12pt]{ctexart}

\usepackage[a4paper,margin=2.35cm]{geometry}
\usepackage{amsmath,amssymb,bm}
\usepackage{booktabs,longtable,tabularx,array}
\usepackage{pdflscape}
\usepackage{enumitem}
\usepackage{xcolor}
\usepackage{microtype}
\usepackage{hyperref}
\usepackage{url}

\hypersetup{
  colorlinks=true,
  linkcolor=blue!55!black,
  citecolor=blue!55!black,
  urlcolor=blue!55!black,
  pdftitle={CISPO 补充材料 S6：综合需求侧灵活性},
  pdfauthor={National\_model project}
}

\setlength{\parindent}{2em}
\setlength{\parskip}{0.25em}
\setlist{nosep,leftmargin=2.2em}
\renewcommand{\arraystretch}{1.18}
\newcolumntype{Y}{>{\raggedright\arraybackslash}X}
\newcommand{\code}[1]{\nolinkurl{#1}}
\newcommand{\unit}[1]{\,\mathrm{#1}}
\newcommand{\casebase}{\code{base}}
\newcommand{\casev}{\code{flex_integrated_v5_central}}
\newcommand{\status}[1]{\nolinkurl{#1}}

\renewcommand{\thesection}{S6.\arabic{section}}
\renewcommand{\thesubsection}{S6.\arabic{section}.\arabic{subsection}}
\renewcommand{\theequation}{S6.\arabic{equation}}
\renewcommand{\thetable}{S6.\arabic{table}}
\renewcommand{\thefigure}{S6.\arabic{figure}}

\title{\textbf{补充材料 S6：综合需求侧灵活性负荷的重建、物理约束、成本与可靠容量价值}}
\author{CISPO National\_model}
\date{2026 年 7 月 30 日方法版本}

\begin{document}
\maketitle

\begin{abstract}
本节给出 CISPO 全国电力系统规划模型中综合需求侧灵活性情景
\casev{} 的完整方法合同。该情景在不改变 Base 年度需求、负荷分解、供给技术、
网络、安全约束和成本边界的前提下，将供暖、制冷、有序充电（V1G）、
双向车网互动（V2G）以及经降额的可靠容量信用放入同一个线性规划。
冷热与电动汽车均以显式服务状态约束能量搬移；V1G、V2G 和建筑侧控制均支付
非零真实资源成本；容量充裕度仍以全年不可变基线峰值为分母，只允许在包含该
峰值的连续四小时窗口内可交付的灵活服务获得降额信用。因此，本情景不是把
``effective peak'' 直接替换为一个可被模型任意压低的峰值，也不把补贴和电价
转移重复计入社会成本。中央参数均登记了低值和高值；其中参与率、服务持续时间
和容量降额系数属于待敏感性检验的工程假设，而非中国实证 ELCC。所有不足
8760 小时的运行均仅用于工程门禁，不能解释为年度科学结果。
\end{abstract}

\noindent\textbf{状态说明。}
本稿描述当前 V5 模型与输入合同，不报告仍在门禁过程中的数值结果。
只有满足 \status{OPTIMAL}、\status{solution_qc=PASS}、全部 hard checks、
当前 input manifest 和有效 result manifest 的 8760 小时结果，才可进入正文
年度比较。

\section{研究问题、情景身份与不可变边界}

\subsection{只保留一个 Base 和一个综合灵活性主情景}

正式主分析集合严格包含两个情景：
\begin{enumerate}
  \item \casebase：需求不可调，所有供暖、制冷和 EV 充电均按外生逐小时曲线
  进入电力平衡；
  \item \casev：在相同年度需求和相同基础系统上，同时启用冷热服务、
  付费 V1G、内生付费 V2G 与降额可靠容量信用。
\end{enumerate}

二者共享同一 baseline contract
\code{base_2024_vre_wave_on_flex_off_v1}。V5 是 analysis case overlay，
其 \code{parent_baseline_case_id} 指向上述 baseline contract，而不是另一个
科学基线。历史 V3/V4、单独 V2G、单独 effective-peak 等配置仅作为工程演进
记录保留，不进入默认论文情景组。双层身份的目的，是把``共同不变的科学基线''
与``本次反事实分析改变了什么''分开记录，避免非 Base 文件仍携带 Base 标签，
也避免 overlay 意外改写固定科学边界。

\subsection{V5 不改变的内容}

V5 不改变：
\begin{itemize}
  \item 31 个省级区域、规划年份、8760 小时时间轴及省际网络；
  \item 年度需求总量、基础残余负荷、供暖、制冷和 EV 四分量的上游重建；
  \item 风、光、波浪、水电、抽蓄、储能、火电、核电和碳移除技术集合；
  \item 逐小时功率平衡、网络约束、备用、惯量、碳约束和 CCS accounting；
  \item 全部货币量的 2025 年不变价人民币口径。
\end{itemize}

因此，Base 与 V5 的差异仅来自需求服务能否在物理边界内跨时移位、是否允许
有限 V2G 放电、这些服务是否支付成本，以及其在峰值窗口能否提供可靠容量。

\section{从逐小时负荷曲线到可调服务分量}

\subsection{历史总负荷锚点和春节形状校正}

上游负荷模块首先构建 2020--2024 年省级北京时间逐小时历史序列。论文来源的
原始总负荷作为不可变锚点；春节校正仅调整小时形状，并在每个省--年满足
\begin{equation}
  \sum_t L^{\mathrm{spr}}_{p,t}
  =
  \sum_t L^{\mathrm{author}}_{p,t}.
\end{equation}
因此，春节处理不创造或删除年度电量。该处理发生在 V5 之前，V5 不重新拟合
春节系数。

\subsection{BAIT 气象暴露与冷热分解}

供暖和制冷分量由省内用电权重聚合的 ERA5 气象暴露以及 BAIT 结构得到
（英文全称为 building-adjusted indoor temperature）。先计算室外体感指标
\begin{align}
  \mathrm{oBAIT}_{p,t}
  &=
  T_{p,t}
  +0.012\!\left(S_{p,t}-k_{S,p,t}\right)
  -0.2\!\left(W_{p,t}-k_{W,p,t}\right) \nonumber\\
  &\quad
  +0.05\!\left(H_{p,t}-k_{H,p,t}\right)
  \operatorname{sign}(T_{p,t}-16),\\
  k_{S,p,t}&=100+7T_{p,t},\\
  k_{W,p,t}&=4.5-0.025T_{p,t},\\
  k_{H,p,t}&=\exp(1+0.06T_{p,t}),
\end{align}
再经 48 小时指数记忆和室内--室外温度混合得到 BAIT。集中供暖区和非集中供暖区
采用不同的 HDD/CDD 阈值。供暖和制冷是同一总负荷的内部拆分，并非叠加到
总负荷上的新增需求。

V5 的冷热物理包络来自同一 BAIT 模型在温度舒适边界
$\pm1\,^{\circ}\mathrm{C}$ 下的逐省逐时差分。这样，允许的上调和下调首先由
气象、原始负荷和舒适边界确定，再乘以显式签约比例；它不是按系统需要自由生成
的储能容量。建筑热惯性和舒适状态采用等效储能表述，相关文献表明该表述适合
刻画建筑负荷的功率、能量和持续时间边界
\cite{IEAEBC67,Stinner2016,Klein2017,HuangWu2019,Chen2023Cooling}。

\subsection{EV 外生充电曲线与未来投影}

历史 EV 分量由省级新能源汽车保有量、月度单车日充电电量和归一化的
96 点充电行为概率生成。未来情景采用省级中情景 NEV 保有量：
\begin{equation}
  L^{\mathrm{EV},0}_{p,h,y}
  =
  \frac{
    N^{\mathrm{EV}}_{p,y}\,
    e_{m(h)}\,
    \pi_h
  }{1000},
\end{equation}
其中 $N^{\mathrm{EV}}_{p,y}$ 为车辆数，$e_{m(h)}$ 为月度单车日充电电量
（kWh vehicle$^{-1}$ day$^{-1}$），$\pi_h$ 为对应小时行为权重。
V5 不修改车辆渗透率、年度交通电量或原始充电曲线，只把其中一个显式比例分为
可被聚合商管理的服务。

\subsection{四分量年度闭合}

不可变基线负荷分解为
\begin{equation}
  L^0_{p,t}
  =
  L^{\mathrm{res}}_{p,t}
  +L^{\mathrm{heat},0}_{p,t}
  +L^{\mathrm{cool},0}_{p,t}
  +L^{\mathrm{EV},0}_{p,t}.
  \label{eq:baseline}
\end{equation}
未来基础残余负荷的年度电量由目标总电量扣除冷热与 EV 年电量后得到，再按
历史 8760 小时模板分配。若残余年度电量为负，数据构建过程 fail closed。
V5 只读取已经闭合的四分量；模型中的任何灵活调度都不得反向改变式
\eqref{eq:baseline} 的输入或 manifest。

\section{供暖和制冷：有签约成本的服务状态模型}

\subsection{签约容量与逐时包络}

令 $s\in\{\mathrm{heat},\mathrm{cool}\}$，$K^{s}_p$ 为省级内生签约功率
（GW），$a^s_{p,t}$ 为逐时可用率，$\overline P^{s,+}_{p,t}$ 和
$\overline P^{s,-}_{p,t}$ 为 BAIT $\pm1\,^{\circ}\mathrm{C}$ 包络乘以
签约参与比例后的上调、下调上限。模型约束
\begin{align}
  0\le P^{s,+}_{p,t}
  &\le
  \min\!\left\{\overline P^{s,+}_{p,t},
  a^s_{p,t}K^s_p\right\},\\
  0\le P^{s,-}_{p,t}
  &\le
  \min\!\left\{\overline P^{s,-}_{p,t},
  a^s_{p,t}K^s_p\right\}.
\end{align}
其中 $P^{s,+}$ 表示为了预热或预冷而增加的电负荷，$P^{s,-}$ 表示相对基线
减少的电负荷。中央签约比例为供暖 25\%、制冷 20\%；10\%--40\% 的低高值
用于敏感性分析。国家需求侧管理制度、地方空调负荷调节目标与需求响应实施规则
证明该类可签约资源具有现实制度基础，但并不直接给出全国统一参与率
\cite{NDRCDemand2023,ShanghaiHVAC2025,AnhuiDR2025,ZhejiangHVAC2025}。

\subsection{非负等效服务库存与周期边界}

建筑服务状态 $S^s_{p,t}$（GWh$_\mathrm{service}$）满足
\begin{equation}
  S^s_{p,t}
  =
  \rho^s_p S^s_{p,t-1}
  +\eta^{s,+}_p P^{s,+}_{p,t}
  -\frac{P^{s,-}_{p,t}}{\eta^{s,-}_p},
  \label{eq:thermalstate}
\end{equation}
\begin{equation}
  0\le S^s_{p,t}\le d^s_p K^s_p.
\end{equation}
中央保持率为供暖 0.94 h$^{-1}$、制冷 0.92 h$^{-1}$；持续时间分别为
4 h 和 5 h。$\eta^{s,+}=\eta^{s,-}=1$ 是避免与保持率重复计损的模型约定。

状态在所选求解时域首尾周期闭合。该边界能防止在末小时无偿积累舒适债务，
但 24 h、168 h 或 744 h 的周期状态并不等价于全年季节边界，故全部短时运行
必须标记 \status{TEST_ONLY_TRUNCATED_HORIZON}。逐时实际冷热负荷为
\begin{equation}
  \widetilde L^s_{p,t}
  =
  L^{s,0}_{p,t}+P^{s,+}_{p,t}-P^{s,-}_{p,t}.
\end{equation}

\section{V1G：有序充电不是免费灵活性}

\subsection{固定与可调充电份额}

中央情景将不可变 EV 充电基线的 $\alpha^{\mathrm{V1G}}=0.15$ 纳入有序充电：
\begin{align}
  L^{\mathrm{EV,flex},0}_{p,t}
  &=\alpha^{\mathrm{V1G}}L^{\mathrm{EV},0}_{p,t},\\
  L^{\mathrm{EV,fixed}}_{p,t}
  &=(1-\alpha^{\mathrm{V1G}})L^{\mathrm{EV},0}_{p,t}.
\end{align}
低值和高值为 0.10 与 0.25。北京有序充电试验显示，激励可以改变参与和充电
电量；价格补贴研究也表明驾驶者反应具有显著异质性
\cite{Liu2021OrderlyCharging,Zhao2026Charging}。因此 15\% 是由现实量级约束的
中央情景假设，而不是全国未来参与率的点预测。

\subsection{移动服务状态与功率约束}

令 $P^{\mathrm{ch}}_{p,t}$ 为被管理车队从电网充电的功率，
$D^{\mathrm{mob}}_{p,t}$ 为外生驾驶能量提取，$E^{\mathrm{EV}}_{p,t}$ 为聚合
移动服务状态。V1G 与 V2G 共用一个物理状态：
\begin{equation}
  E^{\mathrm{EV}}_{p,t}
  =
  E^{\mathrm{EV}}_{p,t-1}
  +\eta^{\mathrm{ch}}P^{\mathrm{ch}}_{p,t}
  -D^{\mathrm{mob}}_{p,t}
  -\frac{P^{\mathrm{V2G,out}}_{p,t}}{\eta^{\mathrm{dis}}}.
  \label{eq:evstate}
\end{equation}
同时满足逐时车队能量上限、最低出发能量和周期首尾闭合。充电功率受外生
逐时可用包络和内生 V1G 签约容量 $K^{\mathrm{V1G}}_p$ 共同约束：
\begin{equation}
  0\le P^{\mathrm{ch}}_{p,t}
  \le
  \min\left\{
    \overline P^{\mathrm{ch}}_{p,t},
    a^{\mathrm{EV}}_{p,t}K^{\mathrm{V1G}}_p
  \right\}.
\end{equation}
输入中的 \code{connected_vehicle_fraction=1} 仅是``已签约服务的可用率归一化''，
不是对全国车辆逐时插枪概率的实证估计。当前上游数据没有逐车出行链、接入会话
或出发 SOC，因此使用聚合服务包络，并将功率比、状态时长和参与率均登记为
待敏感性参数。

\subsection{只对原始充电的向下搬移计费一次}

有序充电不免费。签约容量支付年度控制、计量、聚合和可用性成本。
此外，引入非负变量
\begin{equation}
  R^{\mathrm{V1G}}_{p,t}
  \ge
  L^{\mathrm{EV,flex},0}_{p,t}
  -P^{\mathrm{ch}}_{p,t},
  \qquad
  R^{\mathrm{V1G}}_{p,t}\ge0.
  \label{eq:v1grelocation}
\end{equation}
只有把原始无序充电从参考小时移走的电量才支付 V1G activation/inconvenience
成本。由于 V2G 放电会引起额外回充，若直接使用
$|P^{\mathrm{ch}}-L^{\mathrm{EV,flex},0}|$，会把 V2G 补能误计为 V1G 搬移；
式 \eqref{eq:v1grelocation} 避免这一双重计费。中央年度签约成本为
60 CNY kW$^{-1}$ y$^{-1}$，搬移成本为 300 CNY MWh$^{-1}$。

\section{V2G：嵌套在 V1G 内的内生付费合同}

\subsection{合同嵌套与全国规模上限}

V2G 不是对全部 EV 车队无偿调用的电池。省级内生 V2G 合同功率满足
\begin{equation}
  0\le K^{\mathrm{V2G}}_p
  \le K^{\mathrm{V1G}}_p,
  \qquad
  \sum_p K^{\mathrm{V2G}}_p
  \le \overline K^{\mathrm{V2G}}_y.
  \label{eq:v2gcontract}
\end{equation}
中央全国上限为 2030/2040/2050/2060 年的 10/20/30/40 GW。中国车网互动
政策提出到 2030 年形成千万千瓦级双向互动能力，2025 年首批试点覆盖
9 个城市和 30 个项目
\cite{NDRCVGI2023,NDRCVGIPilot2025}。因此 2030 年 10 GW 仅作为与政策量级
一致的情景上限，不是法律容量上限；2040--2060 年数值是显式外推，必须与
低高情景共同解释。

\subsection{放电、能量和出行边界}

V2G 放电满足
\begin{equation}
  0\le P^{\mathrm{V2G,out}}_{p,t}
  \le
  \min\left\{
    \overline P^{\mathrm{V2G}}_{p,t},
    a^{\mathrm{EV}}_{p,t}K^{\mathrm{V2G}}_p
  \right\},
\end{equation}
并通过式 \eqref{eq:evstate} 同时受到车队状态上限、最低出发能量、充放电效率
和驾驶提取约束。中央充电与放电效率均为 0.94。中国城市真实出行研究表明，
V2G 价值与停留时间、接入条件、用户意愿、补偿和电池退化同时相关，而不能只用
静态电池容量代表
\cite{Li2025V2G}。本模型用聚合线性包络保留这些约束的最低必要结构，但不宣称
复现逐车行为。

\subsection{V2G 四类真实资源成本}

V2G 支付四类成本：
\begin{enumerate}
  \item 增量年度可用性合同；
  \item 双向充电桩、通信平台和维护的年化基础设施成本；
  \item 按实际放电量支付的车主参与/不便成本；
  \item 按实际放电量计提的电池退化成本。
\end{enumerate}
中央参数分别为 60 CNY kW$^{-1}$ y$^{-1}$、
40 CNY kW$^{-1}$ y$^{-1}$、150 CNY MWh$^{-1}$ 和
400 CNY MWh$^{-1}$。中国 V2G 激励研究给出了不同场景下的参与补偿量级，
电池退化研究则说明循环损耗必须独立核算
\cite{Wan2026V2G,Sagaria2025V2G,IEAV2G2026}。

零售分时电价和政府补贴属于消费者、聚合商和财政之间的转移支付，不额外加入
社会系统成本；控制设备、平台、参与不便和电池寿命损耗属于真实资源成本，必须
进入目标函数。回充电量及其损耗通过逐小时电力平衡和式 \eqref{eq:evstate}
自然承担发电成本，不能再次作为独立购电费用重复计费。

\section{有效负荷、逐时安全约束和容量充裕度}

\subsection{有效负荷只用于运行约束}

V5 的实际电网负荷为
\begin{align}
  L^{\mathrm{eff}}_{p,t}
  &=
  L^{\mathrm{res}}_{p,t}
  +\widetilde L^{\mathrm{heat}}_{p,t}
  +\widetilde L^{\mathrm{cool}}_{p,t} \nonumber\\
  &\quad
  +L^{\mathrm{EV,fixed}}_{p,t}
  +P^{\mathrm{ch}}_{p,t}
  -P^{\mathrm{V2G,out}}_{p,t}.
  \label{eq:effective}
\end{align}
逐小时功率平衡、上/下备用需求和惯量相关的运行负荷均使用
$L^{\mathrm{eff}}$，并强制其非负。需求侧灵活性不获得备用容量信用，避免同一
可调功率同时被容量充裕度和备用约束重复占用。

\subsection{为何不直接使用 endogenous effective peak}

若把容量充裕度分母从 $\max_t L^0_{p,t}$ 改成
$\max_t L^{\mathrm{eff}}_{p,t}$，模型可能通过跨小时重排负荷减少装机需求。
这确实是需求灵活性潜在价值的一部分，但若缺少连续事件时长、逐时可用率和
服务恢复约束，单个峰值变量会给出过高容量信用。V5 因此采用保守的
``baseline peak + physically bounded firm credit''：
\begin{equation}
  C^{\mathrm{credit}}_p
  +\sum_s F^{s,\mathrm{firm}}_p
  \ge
  (1+m)\max_t L^0_{p,t},
  \label{eq:adequacy}
\end{equation}
其中 $C^{\mathrm{credit}}_p$ 是按既有容量信用系数计算的供给侧容量，
$m$ 为容量裕度，$F^{s,\mathrm{firm}}_p$ 是灵活服务的可靠容量信用。
这样既保留需求侧减少发电装机的通道，又不允许它仅凭压低单小时峰值获得无限
容量价值。

\subsection{连续四小时峰值窗口的物理降额}

对每个省，令
\begin{equation}
  t^\star_p=\arg\max_t L^0_{p,t},\qquad
  \mathcal W_p=\{t^\star_p-1,t^\star_p,t^\star_p+1,t^\star_p+2\},
\end{equation}
并按全年时间轴循环处理边界。冷热容量信用受签约容量、四小时窗口内最小下调
包络和最小可用率约束；V1G 受签约容量和窗口内最小可削减基线充电约束；
V2G 还受窗口内最小放电包络及能量/4 h 约束。可写为
\begin{align}
 F^{s,\mathrm{firm}}_p
 &\le
 \delta_s\min\left\{
 K^s_p,
 \min_{t\in\mathcal W_p}\overline P^{s,-}_{p,t}
 \right\},
 \quad s\in\{\mathrm{heat},\mathrm{cool}\},\\
 F^{\mathrm{V1G},\mathrm{firm}}_p
 &\le
 \delta_{\mathrm{V1G}}\min\left\{
 K^{\mathrm{V1G}}_p,
 \min_{t\in\mathcal W_p}L^{\mathrm{EV,flex},0}_{p,t}
 \right\},\\
 F^{\mathrm{V2G},\mathrm{firm}}_p
 &\le
 \delta_{\mathrm{V2G}}\min\left\{
 K^{\mathrm{V2G}}_p,
 \min_{t\in\mathcal W_p}\overline P^{\mathrm{V2G}}_{p,t},
 \frac{\min_{t\in\mathcal W_p}\overline E^{\mathrm{fleet}}_{p,t}}{4}
 \right\}.
 \label{eq:firm}
\end{align}
实际实现还将冷热信用乘以窗口内最小可用率与内生签约容量。

中央降额系数为冷热 0.60、V1G/V2G 0.50，低高值分别为
0.30--0.80 和 0.20--0.70。资源充裕度文献强调容量认证应反映事件期间可交付性、
技术之间的相互依赖和系统风险，而不能直接采用铭牌功率
\cite{NREL2024RA,ESIG2023Capacity,Mantegna2024RA}。中国规划文献支持将需求响应
和储能容量价值纳入扩张规划
\cite{Huang2023DRCapacity}，但尚不足以校准全国省级逐时 ELCC。因此上述系数
明确标为 \status{ENGINEERING_ASSUMPTION_NOT_CHINA_ELCC}；监管机构的 ELCC
处理仅作为外部方法参照
\cite{CPUC2026IRP}。

\section{目标函数与成本 accounting scope}

\subsection{社会成本边界}

V5 在既有系统总成本上增加
\begin{align}
 C^{\mathrm{flex}}
 &=
 \sum_{p,s}c^{\mathrm{ena}}_s K^s_p
 +10^{-3}\sum_{p,t,s\in\{\mathrm{heat},\mathrm{cool}\}}
 c^{\mathrm{act}}_s
 \left(P^{s,+}_{p,t}+P^{s,-}_{p,t}\right) \nonumber\\
 &\quad
 +10^{-3}\sum_{p,t}c^{\mathrm{reloc}}_{\mathrm{V1G}}
 R^{\mathrm{V1G}}_{p,t}
 +\sum_p c^{\mathrm{infra}}_{\mathrm{V2G}}K^{\mathrm{V2G}}_p \nonumber\\
 &\quad
 +10^{-3}\sum_{p,t}
 \left(c^{\mathrm{owner}}_{\mathrm{V2G}}
 c^{\mathrm{deg}}_{\mathrm{V2G}}\right)
 P^{\mathrm{V2G,out}}_{p,t}.
 \label{eq:flexcost}
\end{align}
当功率单位为 GW、逐时能量为 GWh 时，年度 CNY/kW 成本乘以 GW，以及
$10^{-3}$ 乘以 CNY/MWh 变动成本，均得到模型内部的 million CNY。
IEA 对需求灵活性的综述同样强调，自动化、聚合、市场接入和用户激励具有非零
实施成本
\cite{IEADemandFlex2025}。

\subsection{避免免费容量和重复计费}

成本与容量约束共同关闭四条潜在``钻空子''通道：
\begin{enumerate}
  \item 无法无限签约：所有服务受逐省外生功率/能量包络约束；
  \item 无法免费获得灵活性：冷热、V1G、V2G 均有年度成本，实际激活另行计费；
  \item 无法把 V2G 补能重复算成 V1G 搬移：只对式
  \eqref{eq:v1grelocation} 的向下位移计费；
  \item 无法仅压低一个小时就减少装机：容量信用必须在连续四小时窗口可交付并
  经过降额。
\end{enumerate}

\section{中央参数、低高区间与证据等级}

\begin{landscape}
\small
\begin{longtable}{
  >{\raggedright\arraybackslash}p{5.2cm}
  >{\centering\arraybackslash}p{2.0cm}
  >{\centering\arraybackslash}p{2.8cm}
  >{\raggedright\arraybackslash}p{3.2cm}
  >{\raggedright\arraybackslash}p{7.4cm}}
\caption{V5 中央参数、敏感性区间与解释边界。完整机器可读值见
\code{config/flexible_load_v5_central_parameters.csv}。}
\label{tab:parameters}\\
\toprule
参数 & 中央值 & 低值--高值 & 证据等级 & 用途与解释边界\\
\midrule
\endfirsthead
\multicolumn{5}{l}{\small 表 \thetable{}（续）}\\
\toprule
参数 & 中央值 & 低值--高值 & 证据等级 & 用途与解释边界\\
\midrule
\endhead
\midrule
\multicolumn{5}{r}{续下页}\\
\endfoot
\bottomrule
\endlastfoot
供暖签约比例 & 0.25 & 0.10--0.40 &
\status{SCENARIO_ASSUMPTION_REQUIRES_SENSITIVITY} &
缩放 BAIT $\pm1\,^{\circ}$C 供暖包络，不改变供暖基线\\
制冷签约比例 & 0.20 & 0.10--0.40 &
\status{SCENARIO_ASSUMPTION_REQUIRES_SENSITIVITY} &
缩放 BAIT $\pm1\,^{\circ}$C 制冷包络，不改变制冷基线\\
供暖/制冷保持率 & 0.94/0.92 & 0.90--0.98/0.85--0.97 &
\status{SCENARIO_ASSUMPTION_REQUIRES_SENSITIVITY} &
等效服务库存逐小时保持率\\
供暖/制冷持续时间 & 4/5 h & 2--8/2--10 h &
\status{SCENARIO_ASSUMPTION_REQUIRES_SENSITIVITY} &
服务能量上限与签约功率的比值\\
冷热年度 enablement cost & 40 & 10--120 CNY/kW-y &
\status{ENGINEERING_ASSUMPTION} &
控制、计量、聚合和可用性真实资源成本\\
冷热 activation cost & 300 & 100--1000 CNY/MWh &
\status{ENGINEERING_ASSUMPTION} &
服务调用和舒适不便成本；地方补贴仅作量级参照\\
V1G 参与率 & 0.15 & 0.10--0.25 &
\status{SCENARIO_ASSUMPTION_REQUIRES_SENSITIVITY} &
不可变 EV 基线中纳入有序充电的比例\\
V2G 参与率 & 0.10 & 0.05--0.30 &
\status{SCENARIO_ASSUMPTION_REQUIRES_SENSITIVITY} &
构造省级 V2G 逐时功率包络；不等于逐时插枪概率\\
EV 充/放电效率 & 0.94/0.94 & 0.90--0.97 &
\status{SCENARIO_ASSUMPTION_REQUIRES_SENSITIVITY} &
电网与聚合移动服务状态之间的能量损耗\\
V1G 充电功率/日均功率比 & 2.0 & 1.5--3.0 &
\status{SCENARIO_ASSUMPTION_REQUIRES_SENSITIVITY} &
有序充电功率上限\\
EV 服务库存持续时间 & 24 h & 12--36 h &
\status{SCENARIO_ASSUMPTION_REQUIRES_SENSITIVITY} &
聚合车队服务状态上限，不是单车电池容量\\
V1G enablement cost & 60 & 20--180 CNY/kW-y &
\status{ENGINEERING_ASSUMPTION} &
年度智能充电控制、聚合和可用性成本\\
V1G relocation cost & 300 & 100--400 CNY/MWh &
\status{ENGINEERING_ASSUMPTION} &
仅对原始充电向下搬移计费一次\\
V2G availability cost & 60 & 20--180 CNY/kW-y &
\status{ENGINEERING_ASSUMPTION} &
在 V1G 合同之外的增量 V2G 可用性成本\\
V2G infrastructure cost & 40 & 20--100 CNY/kW-y &
\status{ENGINEERING_DERIVATION} &
双向充电桩、平台及维护的年化真实资源成本\\
V2G owner compensation & 150 & 100--300 CNY/MWh &
\status{ENGINEERING_ASSUMPTION} &
按放电量计算的参与不便成本，不含退化\\
V2G degradation cost & 400 & 250--500 CNY/MWh &
\status{ENGINEERING_DERIVATION} &
按放电量计算的电池退化成本\\
冷热 firm derating & 0.60 & 0.30--0.80 &
\status{ENGINEERING_ASSUMPTION_NOT_CHINA_ELCC} &
四小时峰值窗口可交付功率的降额；不是中国 ELCC\\
V1G/V2G firm derating & 0.50/0.50 & 0.20--0.70 &
\status{ENGINEERING_ASSUMPTION_NOT_CHINA_ELCC} &
四小时峰值窗口充电削减/放电的降额\\
firm event duration & 4 h & 2--6 h &
\status{ENGINEERING_ASSUMPTION} &
决定容量认证窗口和 V2G 能量/功率约束\\
V2G 全国上限（2030） & 10 GW & 5--10 GW &
\status{POLICY_SCALE_ANCHORED_SCENARIO_CAP} &
与 2030 年千万千瓦级政策目标量级一致，不是法定上限\\
V2G 全国上限（2040/2050/2060） & 20/30/40 GW &
10--30/15--50/20--70 GW &
\status{SCENARIO_EXTRAPOLATION} &
远期外推，必须做敏感性解释\\
\end{longtable}
\end{landscape}

参数证据分为四类：（i）可复现的本地上游曲线和 manifest；
（ii）中国政策、试点和同行评审文献；（iii）国际官方/技术方法；
（iv）明确登记的工程推导。来源注册表逐个解析 parameter source ID，并分别
计算独立来源数、中国证据数和同行评审证据数。``来源存在''不等于``参数已经
校准''；参与率、成本和降额仍必须通过低/中央/高敏感性报告。

\section{输入生成、manifest 和结果审计}

\subsection{确定性输入生成}

V5 小时输入包括：
\begin{itemize}
  \item 省--年--小时冷热上/下调包络；
  \item 省级冷热保持率、效率和持续时间；
  \item EV 逐时充电/放电功率包络、车队服务能量上限和连接归一化；
  \item 外生驾驶能量提取与最低出发能量；
  \item 四类服务的年度和逐能量成本。
\end{itemize}
构建器 \code{scripts/build_flexible_load_v5_inputs.py} 从已接受的负荷分解、
V3 舒适包络和中央参数表确定性生成上述文件，并写入 SHA256 manifest。
验证器 \code{scripts/validate_flexible_load_v5_inputs.py} 不信任手写 PASS：
它重新解析每个 source ID、重算证据计数、核对四个规划年、参数顺序、
省份--年份--小时完整性、能量闭合和 manifest 哈希，任何不一致均 fail closed。

\subsection{优化结果必须输出的 V5 专用证据}

除既有供给、网络、备用、惯量、碳、CCS 和成本结果外，V5 至少输出：
\begin{itemize}
  \item 基线与有效负荷、冷热实际负荷和服务状态；
  \item V1G 签约功率、实际充电、向下搬移量和车队服务状态；
  \item V2G 签约功率、放电量、全国上限和合同嵌套残差；
  \item 四类 firm capacity credit 及各自物理上界；
  \item enablement、thermal activation、V1G relocation、V2G
  infrastructure、owner participation 和 degradation 成本分项；
  \item 状态转移、周期边界、有效负荷重建、同时充放电、容量信用上界等 hard
  checks。
\end{itemize}

\subsection{分级工程门禁}

代码与输入合同完成后依次执行：
\begin{enumerate}
  \item 单元测试、输入审计和 release-contract 审计；
  \item 新输出根下的 1 h Base/V5 结构门禁；
  \item 新输出根下的 24 h Base/V5 机制门禁；
  \item 2030--2060 四年 168 h Base/V5 planning sequence；
  \item 经明确授权的单个 2030/744 h cold engineering gate。
\end{enumerate}
所有运行必须串行，不允许第二个并发求解、basis gate 或已拒绝的
\code{Crossover=3}。744 h 仍为 \status{TEST_ONLY_TRUNCATED_HORIZON}，
不能作为年度结果或 2040 state anchor。通过 744 h 也不自动启动 8760 h。

\section{结果解释框架与限制}

\subsection{Base--V5 比较应报告什么}

完整年度结果应至少同时报告：
\begin{itemize}
  \item 总系统成本及 V5 六项增量成本；
  \item 发电、储能、输电和 V2G 合同容量变化；
  \item 基线峰值、有效峰值及四类 firm capacity credit；
  \item 逐小时负荷转移、V2G 放电/回充损耗和未满足服务检查；
  \item 碳约束、备用、惯量、网络和水电运行是否仍满足；
  \item 中央值相对低/高参与率、成本和 derating 的敏感性。
\end{itemize}
``容量减少''只有在式 \eqref{eq:adequacy} 与式 \eqref{eq:firm} 同时通过、
成本分项完整、且 Base/V5 使用相同输入 manifest 时，才能解释为灵活性资源的
可靠容量价值。仅观察有效峰值下降不能证明可替代发电装机。

\subsection{当前限制}

\begin{enumerate}
  \item 建筑舒适包络来自聚合 BAIT $\pm1\,^{\circ}$C，不是逐建筑热工模拟；
  \item EV 使用省级聚合服务状态，未显式表示逐车出行链、接入会话、充电站
  配电容量或单车 SOC；
  \item 远期 V2G 全国上限和全部 derating 系数是情景参数，不是预测；
  \item 四小时窗口是透明的工程容量认证近似，不替代概率型 ELCC/LOLE 计算；
  \item 周期时域边界适合工程门禁和全年闭合，但短时门禁会人为连接时域首尾；
  \item 需求侧资源当前不提供备用信用，可能低估其部分运行价值，但能避免容量
  与备用双重计数；
  \item 只有 8760 h accepted Base 与 V5 才能支撑论文年度结论；短时求解只能
  排查代码、数据、数值稳定性和 accounting 漏洞。
\end{enumerate}

\clearpage
\begin{thebibliography}{99}

\bibitem{IEAEBC67}
IEA Energy in Buildings and Communities Programme.
\newblock Annex 67: Energy Flexible Buildings.
\newblock International Energy Agency, 2019.
\newblock \url{https://www.iea-ebc.org/projects/project?AnnexID=67}.

\bibitem{Stinner2016}
Stinner S, Huchtemann K, Müller D.
\newblock Quantifying the operational flexibility of building energy systems
with thermal energy storages.
\newblock \emph{Applied Energy}, 181:140--154, 2016.
\newblock \url{https://doi.org/10.1016/j.apenergy.2016.08.055}.

\bibitem{Klein2017}
Klein K, Herkel S, Henning H-M, Felsmann C.
\newblock Load shifting using the heating and cooling system of an office
building: quantitative potential evaluation for different flexibility and
storage options.
\newblock \emph{Applied Energy}, 203:917--937, 2017.
\newblock \url{https://doi.org/10.1016/j.apenergy.2017.06.073}.

\bibitem{HuangWu2019}
Huang S, Wu D.
\newblock Validation on aggregate flexibility from residential air
conditioning systems for building-to-grid integration.
\newblock \emph{Energy and Buildings}, 200:58--67, 2019.
\newblock \url{https://doi.org/10.1016/j.enbuild.2019.07.043}.

\bibitem{Chen2023Cooling}
Chen Q, Kuang Z, Liu X, Zhang T.
\newblock Techno-economic comparison of cooling storage and battery for
electricity flexibility.
\newblock \emph{Energy Conversion and Management}, 289:117183, 2023.
\newblock \url{https://doi.org/10.1016/j.enconman.2023.117183}.

\bibitem{NDRCDemand2023}
国家发展和改革委员会.
\newblock 《电力需求侧管理办法（2023 年版）》, 2023.
\newblock \url{https://www.ndrc.gov.cn/xxgk/zcfb/ghxwj/202309/t20230927_1360902.html}.

\bibitem{ShanghaiHVAC2025}
上海市人民政府.
\newblock 《关于做好本市 2025 年空调负荷调控能力建设有关工作的通知》,
2025.
\newblock \url{https://www.shanghai.gov.cn/qtsj/20250617/9db446ddad1744b6af10b6f7205d8aab.html}.

\bibitem{AnhuiDR2025}
安徽省能源局.
\newblock 《安徽省电力需求响应实施细则》, 2025.
\newblock \url{https://www.ahjd.gov.cn/Jczwgk/show/3553986.html}.

\bibitem{ZhejiangHVAC2025}
衢州市人民政府.
\newblock 《大科创专项政策实施细则（2025 年版）》, 2025.
\newblock \url{https://zjjcmspublic.oss-cn-hangzhou-zwynet-d01-a.internet.cloud.zj.gov.cn/jcms_files/jcms1/web3084/site/attach/0/adde479fede6497d921889690afad168.pdf}.

\bibitem{IEADemandFlex2025}
International Energy Agency.
\newblock \emph{The Value of Demand Flexibility}.
\newblock IEA, Paris, 2025.
\newblock \url{https://www.iea.org/reports/the-value-of-demand-flexibility}.

\bibitem{NDRCVGI2023}
国家发展和改革委员会, 国家能源局等.
\newblock 《关于加强新能源汽车与电网融合互动的实施意见》, 2023.
\newblock \url{https://zfxxgk.nea.gov.cn/2023-12/13/c_1310758710.htm}.

\bibitem{NDRCVGIPilot2025}
国家发展和改革委员会, 国家能源局等.
\newblock 《首批车网互动规模化应用试点》, 2025.
\newblock \url{https://www.ndrc.gov.cn/xxgk/zcfb/tz/202504/t20250402_1396958.html}.

\bibitem{Liu2021OrderlyCharging}
Liu X, Li X, Chen H, Xu H.
\newblock Aggregate electric vehicles in demand side for ancillary service.
\newblock In: \emph{34th International Electric Vehicle Symposium and
Exhibition (EVS34)}, Nanjing, 2021.
\newblock \url{https://www.citelec.org/EVS34/download.php?f=papers/evs34-5010511.pdf}.

\bibitem{Zhao2026Charging}
Zhao X, Chen H, Qiu Y, Deng Y.
\newblock Impact of charging price subsidies on the charging behavior of
heterogeneous electric vehicle drivers: evidence from Beijing.
\newblock \emph{Journal of the Association of Environmental and Resource
Economists}, 13(3):713--754, 2026.
\newblock \url{https://doi.org/10.1086/739399}.

\bibitem{IEAV2G2026}
International Energy Agency.
\newblock \emph{Vehicle-to-Grid Technology}.
\newblock IEA, Paris, 2026.
\newblock \url{https://www.iea.org/reports/vehicle-to-grid-technology}.

\bibitem{Li2025V2G}
Li K, Li X, Xiong Z, Tao S, Zhao G, Jiang Y, Qi H, Zhang Y.
\newblock Unlocking vehicle-to-grid potential of load shifting in China's
megacities considering comprehensive real-world behaviors.
\newblock \emph{Nature Communications}, 16:10087, 2025.
\newblock \url{https://doi.org/10.1038/s41467-025-65073-8}.

\bibitem{Wan2026V2G}
Wan M, Xia L, Huo Y, Gong Y, Geng G, Jiang Q.
\newblock Case study of vehicle-to-grid incentive design under multiple
scenarios: a policy-making perspective.
\newblock \emph{IET Generation, Transmission \& Distribution},
20(1):e70254, 2026.
\newblock \url{https://doi.org/10.1049/gtd2.70254}.

\bibitem{Sagaria2025V2G}
Sagaria S, van der Kam M, Boström C.
\newblock Vehicle-to-grid impact on battery degradation and estimation of V2G
economic compensation.
\newblock \emph{Applied Energy}, 377:124546, 2025.
\newblock \url{https://doi.org/10.1016/j.apenergy.2024.124546}.

\bibitem{NREL2024RA}
National Renewable Energy Laboratory.
\newblock \emph{Explained: Fundamentals of Power Grid Reliability and Clean
Electricity}. NREL/TP-6A40-85880, 2024.
\newblock \url{https://docs.nrel.gov/docs/fy24osti/85880.pdf}.

\bibitem{ESIG2023Capacity}
Energy Systems Integration Group.
\newblock \emph{Ensuring Efficient Reliability: New Design Principles for
Capacity Accreditation}, 2023.
\newblock \url{https://www.esig.energy/wp-content/uploads/2023/02/ESIG-Design-principles-capacity-accreditation-report-2023.pdf}.

\bibitem{Mantegna2024RA}
Mantegna G, Huang Z, Van Caelenberg G, Frew B, Lynch M, O'Malley M.
\newblock Maintaining reliability while navigating unprecedented uncertainty:
a synthesis of and guide to advances in electric sector resource adequacy.
\newblock \emph{arXiv:2412.00533}, 2024.
\newblock \url{https://doi.org/10.48550/arXiv.2412.00533}.

\bibitem{Huang2023DRCapacity}
Huang Y, Zhang Y, Xia Z, Wang H, Wu M, Wang N, Chen Q, Zhu T, Chen X.
\newblock Power system planning considering demand response resources and
capacity value of energy storage.
\newblock \emph{Journal of Shanghai Jiao Tong University},
57(4):432--441, 2023.
\newblock \url{https://doi.org/10.16183/j.cnki.jsjtu.2021.477}.

\bibitem{CPUC2026IRP}
California Public Utilities Commission.
\newblock Inputs and assumptions for 2026 integrated resource planning, 2026.
\newblock \url{https://www.cpuc.ca.gov/industries-and-topics/electrical-energy/electric-power-procurement/long-term-procurement-planning/2024-26-irp-cycle-events-and-materials}.

\end{thebibliography}

\end{document}
